\documentclass[../main]{subfiles}
\ifSubfilesClassLoaded{\setcounter{chapter}{0}}{}
\begin{document}

\chapter{Pendahuluan dan Orientasi Buku}

Buku ajar ini disusun untuk mendukung mata kuliah \textbf{Logika Matematika} (2 SKS, 16 pertemuan) Program Studi Teknik Informatika, Fakultas Teknologi Informasi, Universitas Bale Bandung, berbasis \textbf{Outcome-Based Education (OBE)}. Bab ini memperkenalkan tujuan buku, keterkaitannya dengan RPS, cara menggunakan buku, serta pemetaan CPL, CPMK, dan Sub-CPMK ke bab-bab dalam buku.

\section*{Tujuan Buku Ajar dan Keterkaitan dengan RPS}

Buku ini dirancang agar mahasiswa mencapai capaian pembelajaran mata kuliah (CPMK) dan sub-CPMK yang ditetapkan dalam RPS. Setiap bab unit materi inti (Bab 3--14) mencakup satu atau lebih Sub-CPMK dengan judul bab yang selaras 100\% dengan materi pembelajaran dalam silabus. Bab 15 berisi evaluasi, refleksi, dan integrasi kompetensi (Kuis, UTS, UAS, Tugas) sesuai bobot penilaian RPS.

\section*{Capaian Pembelajaran Lulusan (CPL)}

CPL yang dibebankan pada mata kuliah ini mencakup kompetensi lulusan dalam aspek pengetahuan, keterampilan, dan sikap:

\begin{itemize}[leftmargin=*]
  \item \textbf{CPL01:} Kemampuan menganalisis persoalan computing yang kompleks serta menerapkan prinsip-prinsip computing dan disiplin ilmu relevan lainnya untuk mengidentifikasi solusi, dengan mempertimbangkan wawasan perkembangan ilmu transdisiplin.
  \item \textbf{CPL02:} Memiliki pengetahuan yang memadai terkait dengan cara kerja sistem komputer dan mampu merancang dan mengembangkan berbagai algoritma metode untuk memecahkan masalah.
\end{itemize}

\section*{Capaian Pembelajaran Mata Kuliah (CPMK)}

Kemampuan atau kompetensi spesifik yang diharapkan dicapai setelah menyelesaikan mata kuliah:

\begin{itemize}[leftmargin=*]
  \item \textbf{CPMK1 (Fokus Logika Proposisional \& Analisis Argumen -- Terkait CPL01):} Mahasiswa mampu menganalisis validitas argumen dan persoalan komputasi dasar dengan menerapkan prinsip-prinsip logika proposisional, tabel kebenaran, dan penarikan kesimpulan (silogisme, modus ponens, modus tollens) untuk mengidentifikasi solusi yang logis.
  \item \textbf{CPMK2 (Fokus Logika Predikat \& Pemodelan -- Terkait CPL01 \& CPL02):} Mahasiswa mampu memodelkan persoalan yang lebih kompleks ke dalam bentuk ekspresi logika predikat (menggunakan kuantor universal dan eksistensial) sebagai langkah awal dalam merancang metode penyelesaian masalah yang terstruktur.
  \item \textbf{CPMK3 (Fokus Aljabar Boolean \& Sistem Komputer -- Terkait CPL02):} Mahasiswa mampu menerapkan hukum-hukum aljabar Boolean untuk menyederhanakan fungsi logika dan merepresentasikannya ke dalam bentuk gerbang logika dasar, sebagai bentuk pemahaman terhadap cara kerja perangkat keras sistem komputer.
  \item \textbf{CPMK4 (Fokus Teknik Pembuktian \& Validasi Algoritma -- Terkait CPL01 \& CPL02):} Mahasiswa mampu merancang argumen matematis dan membuktikan kebenaran suatu metode atau algoritma sederhana menggunakan teknik pembuktian formal (seperti deduksi, kontradiksi, atau induksi matematika) untuk memastikan solusi komputasi berjalan sesuai spesifikasi.
\end{itemize}

\section*{Sub-CPMK / Indikator Pencapaian}

Penjabaran CPMK menjadi indikator yang lebih terukur dan dapat diuji (teks persis RPS):

\begin{subcpmk}
  \item \textbf{Sub-CPMK 1.1:} Mahasiswa mampu menjelaskan konsep dasar proposisi, nilai kebenaran, dan kata hubung logika (konjungsi, disjungsi, negasi, implikasi, biimplikasi).
  \item \textbf{Sub-CPMK 1.2:} Mahasiswa mampu menggunakan tabel kebenaran untuk mengevaluasi pernyataan majemuk serta mengidentifikasi tautologi, kontradiksi, dan kontingensi.
  \item \textbf{Sub-CPMK 1.3:} Mahasiswa mampu menganalisis validitas argumen menggunakan aturan inferensi (modus ponens, modus tollens, silogisme) dalam memecahkan persoalan komputasi dasar.
  \item \textbf{Sub-CPMK 2.1:} Mahasiswa mampu menjelaskan konsep dasar logika predikat, fungsi proposisional, dan ruang lingkup pembicaraan (\textit{universe of discourse}).
  \item \textbf{Sub-CPMK 2.2:} Mahasiswa mampu menerjemahkan pernyataan bahasa alami ke dalam ekspresi logika predikat menggunakan kuantor universal dan eksistensial.
  \item \textbf{Sub-CPMK 2.3:} Mahasiswa mampu memodelkan masalah komputasi yang terstruktur dengan merepresentasikan logika predikat sebagai basis perancangan algoritma.
  \item \textbf{Sub-CPMK 3.1:} Mahasiswa mampu mendefinisikan prinsip-prinsip dan hukum-hukum dasar Aljabar Boolean.
  \item \textbf{Sub-CPMK 3.2:} Mahasiswa mampu menyederhanakan fungsi logika menggunakan hukum Aljabar Boolean atau Peta Karnaugh (K-Map).
  \item \textbf{Sub-CPMK 3.3:} Mahasiswa mampu mengonversi dan merepresentasikan fungsi logika yang telah disederhanakan ke dalam bentuk rangkaian gerbang logika dasar pembentuk sistem komputer.
  \item \textbf{Sub-CPMK 4.1:} Mahasiswa mampu menerapkan berbagai metode pembuktian matematis dasar (pembuktian langsung, pembuktian tidak langsung/kontradiksi, dan kontraposisi).
  \item \textbf{Sub-CPMK 4.2:} Mahasiswa mampu mengaplikasikan prinsip Induksi Matematika untuk membuktikan kebenaran sebuah deret atau properti dari algoritma rekursif sederhana.
  \item \textbf{Sub-CPMK 4.3:} Mahasiswa mampu menyusun argumen matematis formal untuk memvalidasi langkah-langkah dalam penyelesaian masalah komputasi.
\end{subcpmk}

\section*{Pemetaan CPL, CPMK, dan Sub-CPMK ke Bab}

Tabel berikut memetakan setiap Sub-CPMK (dan bab unit materi) ke nomor bab dalam buku ini. Bab 1 dan 2 bersifat pengantar; Bab 3--14 merupakan unit materi inti sesuai 13 butir materi RPS; Bab 15 berisi evaluasi, refleksi, dan integrasi kompetensi (Kuis, UTS, UAS, Tugas).

\begin{table}[H]
\centering
\small
\resizebox{\textwidth}{!}{\begin{tabular}{|c|l|c|c|}
\hline
\textbf{Sub-CPMK} & \textbf{Judul Bab / Materi} & \textbf{CPMK} & \textbf{Bab} \\
\hline
1.1 & Konsep Dasar Proposisi dan Kata Hubung Logika & CPMK1 & 3 \\
\hline
1.2 & Tabel Kebenaran, Tautologi, Kontradiksi, dan Kontingensi & CPMK1 & 4 \\
\hline
1.3 & Aturan Inferensi dan Penarikan Kesimpulan & CPMK1 & 5 \\
\hline
2.1 & Konsep Logika Predikat, Fungsi Proposisional, UoD & CPMK2 & 6 \\
\hline
2.2 & Kuantor Universal dan Eksistensial & CPMK2 & 7 \\
\hline
2.3 & Pemodelan Masalah Komputasi berbasis Logika Predikat & CPMK2 & 8 \\
\hline
3.1 & Prinsip dan Hukum-hukum Dasar Aljabar Boolean & CPMK3 & 9 \\
\hline
3.2 & Penyederhanaan Fungsi Logika (Boolean dan K-Map) & CPMK3 & 10 \\
\hline
3.3 & Representasi Fungsi Logika ke Gerbang Logika Dasar & CPMK3 & 11 \\
\hline
4.1 & Metode Pembuktian Langsung dan Kontraposisi & CPMK4 & 12 \\
\hline
4.1 & Pembuktian Tidak Langsung dan Kontradiksi & CPMK4 & 13 \\
\hline
4.2 & Prinsip Induksi Matematika & CPMK4 & 14 \\
\hline
4.3 & Argumen Matematis Formal dan Validasi Algoritma & CPMK4 & 14 \\
\hline
-- & Evaluasi, Refleksi, dan Integrasi Kompetensi (Kuis, UTS, UAS, Tugas) & Semua & 15 \\
\hline
\end{tabular}}
\caption{Pemetaan Sub-CPMK dan CPMK ke Nomor Bab}
\end{table}

% ============================================================
% MATERI POKOK
% ============================================================
\section{Mengapa Teknik Informatika Memuja Logika?}

Selamat datang di titik awal medan tempur Ilmu Komputer (*Computer Science*). Saat Anda membuka lembaran buku ini, Anda mungkin bertanya-tanya: \textit{"Untuk apa seorang mahasiswa Teknik Informatika yang seharusnya belajar mengetik kode aplikasi, malah dipaksa duduk berjibaku mempelajari deretan huruf P, Q, dan operator V (OR), $\land$ (AND) yang terlihat seperti peninggalan filsuf Yunani Kuno?"} \cite{rosen2019,copi2014}

Jawabannya keras dan sederhana: \textbf{Mengetik kode (Coding) bisa dilakukan oleh siapa saja, bahkan oleh anak-anak menggunakan *Drag-and-Drop*. Namun, Mendesain Arsitektur Rekayasa Perangkat Lunak (\textit{Software Engineering}) dan Algoritma Kecerdasan Buatan (\textit{Artificial Intelligence}) yang bebas \textit{Bug} maut, murni mensyaratkan otak yang telah dikalibrasi oleh Logika Matematika} \cite{levin_discrete,stanford_cs103}.

Buku ajar ini ditulis bukan sebagai seonggok dokumen matematis teoretis yang mengantuk. Ini adalah **Manual Keselamatan (\textit{Survival Guide})** Anda untuk diizinkan masuk ke dalam jajaran Insinyur IT berstandar global. Universitas terkemuka dunia seperti MIT (dalam \textit{Mathematics for Computer Science}) dan Stanford (dalam \textit{Mathematical Foundations of Computing}) \cite{mit_6042,stanford_cs103} secara absolut memelopori bahwa tanpa Logika Proposisional dan Predikat, seorang sarjana komputer tak lebih dari sekadar "Tukang Ketik Berijazah".

\subsection{Menyelamatkan Sistem dari Kiamat \textit{Logical Bug}}

Aplikasi kasir, sistem keamanan reaktor nuklir, hingga *Autopilot* pesawat Boeing, semuanya dijalankan oleh deretan jutaan *If-Else Statement* (Kondisi Percabangan). Ketika seorang *Programmer* buta akan hukum \textit{Kontraposisi} atau gagal mengadopsi \textit{De Morgan's Laws} saat meracik filter *Password* raksasa, mesin bukannya *Error* secara sintaks. Mesin akan terus berjalan mulus, **namun** diam-diam membelokkan arah pesawat menabrak gunung atau membocorkan data triliunan saldo Bank. Di ranah komputasi, kecacatan alur pikir tak kasatmata ini bergelar \textbf{\textit{Logic Bug}}---mimpi buruk yang hanya bisa diredam jika otak Anda terbiasa menguliti tabel kebenaran validitas Argumen \cite{forallx,iep_natded}.

\subsection{Tangga Menuju Kecerdasan Buatan (AI)}

Memasuki era abad 21, Anda niscaya bercita-cita membangun *Machine Learning* atau *Expert Systems* yang sanggup mendiagnosis kanker sehebat dokter \cite{saylor_cs202}. AI modern tidak menebak; ia bernalar murni menggunakan struktur \textbf{Logika Predikat Orde Pertama (\textit{First-Order Logic})}. Buku ini akan mengajarkan langkah demi langkah bagaimana menarik Kesimpulan Mutlak (\textit{Inference}) bertenaga *Modus Ponens*, mengawinkan insting manusia purba dengan komputasi biner abad mesin. Menguasai buku ini, berarti Anda menggenggam cetak biru (*blueprint*) sejati bagaimana *Database Relasional (SQL)* dan Mesin *AI* mengambil keputusan \cite{rosen2019}.

\section{Outcome-Based Education (OBE) dan capaian pembelajaran}

Rencana pembelajaran semester (RPS) mata kuliah ini mengacu pada \textbf{Outcome-Based Education (OBE)}: capaian yang diharapkan dijelaskan secara eksplisit, dan penilaian mengarah pada bukti penguasaan luaran (\textit{outcome}), bukan hanya pada kehadiran atau volume catatan \cite{spady1994,obe_springer2024}.

\subsection{Paradigma luaran versus masukan}

Pada pendekatan berorientasi masukan (\textit{input-based}), indikator sukses sering dikaitkan dengan kegiatan formal (misalnya jumlah pertemuan atau lembar ringkasan). Dalam OBE, penekanan dipindahkan ke \textbf{demonstrasi kemampuan} sesuai capaian yang telah ditetapkan \cite{spady1994}: mahasiswa diharapkan menunjukkan, misalnya, kemampuan menerjemahkan persyaratan informal ke ekspresi logika atau menyusun bukti yang valid---sesuai praktik \textit{requirements engineering} dan standar kompetensi di bidang komputasi.

\subsection{Sub-CPMK sebagai target terukur}

Setiap bab unit materi menyatakan tujuan pembelajaran tingkat lebih rinci melalui \textbf{Sub-CPMK} (Sub-Capaian Pembelajaran Mata Kuliah) \cite{in4obe}. Contoh pemetaan dengan struktur buku ini:

\begin{itemize}
  \item \textbf{Bab 3--5} memfokuskan logika proposisional (proposisi, tabel kebenaran, tautologi/kontradiksi/kontingensi) serta aturan inferensi sesuai Sub-CPMK 1.1--1.3.
  \item \textbf{Bab 6--8} memperluas ke logika predikat, kuantifier, dan pemodelan sesuai Sub-CPMK 2.1--2.3.
  \item \textbf{Bab 9--11} membahas aljabar Boolean, penyederhanaan, dan gerbang logika (Sub-CPMK 3.1--3.3).
  \item \textbf{Bab 12--14} mencakup metode bukti, induksi, serta validasi algoritma (Sub-CPMK 4.1--4.3).
\end{itemize}

Latihan dan aktivitas di akhir bab dirancang sebagai latihan terstruktur yang selaras dengan indikator OBE; peran dosen cenderung sebagai \textbf{fasilitator} yang membimbing pencapaian capaian, bukan sekadar penyaji materi satu arah \cite{obe_springer2024,benari_matlogic}.

\section{Cara menggunakan buku ajar ini}

Materi logika bersifat terstruktur dan saling membangun: konsep di bab awal dipakai kembali pada definisi, bukti, dan aplikasi di bab selanjutnya \cite{levin_discrete}. Oleh karena itu disarankan pembelajaran mandiri (\textit{self-directed learning}) dengan membaca berurutan sesuai silabus, mengerjakan latihan, dan menguji pemahaman secara aktif \cite{rosen2019,forallx}.

\subsection{Urutan baca dan ketergantungan konsep}

Konsep dasar seperti negasi, konjungsi, dan implikasi muncul di bab-bab awal materi inti dan digunakan kembali pada inferensi, logika predikat, hingga pembuktian \cite{rosen2019,forallx}. Melompati bab dapat mempersulit pemahaman langkah bukti atau manipulasi kuantifier di bagian kemudian. Disarankan menempuh alur Bab 3--14 sesuai urutan buku kecuali instruktur memberikan panduan khusus.

\subsection{Memverifikasi pemahaman secara mandiri}

Membaca bukti atau tabel kebenaran hendaknya disertai reproduksi mandiri: misalnya, menyusun kembali tabel kebenaran atau langkah deduksi di kertas kerja sebelum membandingkan dengan solusi \cite{proofs_checker,carnap}. Jika terjebak pada suatu langkah, bandingkan dengan pembahasan setelah mencoba secara independen---pendekatan ini memperkuat keterampilan pemecahan masalah yang relevan dalam praktik rekayasa perangkat lunak.

\subsection{Komponen penilaian di setiap bab}

Setiap bab umumnya memuat:
\begin{enumerate}
  \item \textbf{Aktivitas pembelajaran:} menghubungkan definisi dengan contoh atau studi kasus sederhana.
  \item \textbf{Latihan dan refleksi:} soal esai atau analisis untuk menguji pemahaman konsep dan fallasi argumen \cite{copi2014}.
  \item \textbf{Checklist atau asesmen:} indikator kesiapan terhadap capaian yang dipetakan ke Sub-CPMK.
\end{enumerate}

Gunakan komponen tersebut secara konsisten sebagai panduan belajar, bukan hanya sebagai pelengkap teori \cite{levin_discrete,stanford_cs103}.

\section{Alur materi dan peta konsep}

Bab-bab materi inti tersusun secara berjenjang sesuai pemetaan Sub-CPMK (Gambar~\ref{fig:peta-konsep}). Alur ini membantu menempatkan setiap topik dalam konteks keseluruhan kurikulum OBE \cite{openlogic,levin_discrete}.

\begin{figure}[H]
\centering
\begin{tikzpicture}[
  node distance=0.6cm and 1cm,
  box/.style={rectangle, draw=black, rounded corners, fill=blue!10, text width=2.6cm, minimum height=0.85cm, align=center, font=\small},
  root/.style={rectangle, draw=black, rounded corners, fill=green!20, text width=3.4cm, minimum height=0.9cm, align=center, font=\small\bfseries},
  arrow/.style={->, thick, >=stealth}
]
  \node[root] (root) {Logika Matematika\\(alur kurikulum)};
  \node[box, below left=1.2cm and 0.4cm of root] (b3) {Tahap I\\Bab 3--5:\\Proposisi, tabel, inferensi};
  \node[box, below=1.2cm of root] (b4) {Tahap II\\Bab 6--8:\\Predikat dan pemodelan};
  \node[box, below right=1.2cm and 0.4cm of root] (b5) {Tahap III\\Bab 9--11:\\Boolean, K-Map, gerbang};
  \node[box, below=1.8cm of b4] (b6) {Tahap IV\\Bab 12--14:\\Bukti, induksi, validasi};
  \draw[arrow] (root) -- (b3);
  \draw[arrow] (root) -- (b4);
  \draw[arrow] (root) -- (b5);
  \draw[arrow] (b3) -- (b4);
  \draw[arrow] (b4) -- (b5);
  \draw[arrow] (b5) -- (b6);
  \draw[arrow] (b4) -- (b6);
\end{tikzpicture}
\caption{Skema tahapan materi inti dan nomor bab}
\label{fig:peta-konsep}
\end{figure}

\subsection{Tahap I (Bab 3--5): Proposisi dan inferensi}
Bab 3 dan 4 membahas proposisi, operator logika, serta tabel kebenaran dan klasifikasi formula. Bab 5 memperdalam aturan inferensi dan penarikan kesimpulan pada logika proposisional.

\subsection{Tahap II (Bab 6--8): Predikat, kuantifier, dan pemodelan}
Bab 6--8 mengembangkan logika predikat, kuantifier, serta contoh pemodelan masalah komputasi. Konsep ini menjadi landasan untuk menerjemahkan persyaratan informal ke bahasa formal.

\subsection{Tahap III (Bab 9--11): Aljabar Boolean dan rangkaian}
Bab 9--11 membahas hukum aljabar Boolean, penyederhanaan (termasuk K-Map), dan representasi ke gerbang logika sebagai jembatan ke pemahaman struktur rangkaian digital.

\subsection{Tahap IV (Bab 12--14): Pembuktian dan validasi}
Bab 12--14 menempatkan teknik pembuktian, induksi matematika, serta argumen formal untuk validasi algoritma dan verifikasi program sesuai Sub-CPMK 4.x, termasuk konsep prekondisi, postkondisi, dan loop invariant pada bagian validasi.

\section{Peta Konsep Logika Matematika}

Sebagai ringkasan visual dari konteks kurikulum OBE untuk mata kuliah ini, peta konsep berikut berfungsi sebagai orientasi bagi mahasiswa untuk memahami alur keseluruhan materi Logika Matematika dalam buku ini. Alur pembelajaran mengikuti struktur standar yang digunakan oleh Open Logic Project dan Discrete Mathematics Levin: dimulai dari logika proposisional (proposisi, operator, tabel kebenaran), dilanjutkan ke ekuivalensi dan tautologi, kemudian logika predikat dengan kuantifier, metode inferensi dan pembuktian, induksi matematika, dan terakhir pemodelan logika predikat orde pertama \cite{openlogic,levin_discrete}.

Bab III dan IV membahas logika proposisional secara mendalam; Bab V memperluas ke logika predikat; Bab VI dan VII fokus pada aturan inferensi dan metode bukti termasuk induksi; Bab VIII memperkenalkan pemodelan logika predikat orde pertama sebagai aplikasi praktis. Memahami keterkaitan antarbab membantu mahasiswa menempatkan setiap topik dalam konteks keseluruhan dan melihat bagaimana konsep-konsep tersebut saling membangun.

Berikut peta konsep topik utama per bab:

\begin{center}
\begin{tabular}{|c|l|}
\hline
\textbf{Bab} & \textbf{Topik Utama} \\
\hline
II & Landasan Teori: Logika dan Informatika \\
III & Proposisi $\rightarrow$ Operator $\rightarrow$ Tabel Kebenaran \\
IV & Ekuivalensi $\rightarrow$ Tautologi $\rightarrow$ Translasi \\
V & Predikat $\rightarrow$ Kuantifier $\rightarrow$ Negasi \\
VI & Inferensi $\rightarrow$ Bukti Langsung $\rightarrow$ Kontraposisi \\
VII & Induksi Matematika \\
VIII & Pemodelan FOL: Fakta, Aturan, Inferensi \\
\hline
\end{tabular}
\end{center}

Gunakan peta konsep ini sebagai acuan saat mempelajari setiap bab; rujuk kembali tabel di atas untuk mengingat posisi topik yang sedang dipelajari dalam keseluruhan alur. Hal ini membantu menghindari kehilangan konteks dan memperkuat pemahaman koneksi antar konsep.


% ============================================================
% AKTIVITAS PEMBELAJARAN
% ============================================================
\begin{aktivitas}
  \item \textbf{Target capaian akhir semester:} Rumuskan satu indikator keterampilan terkait logika (misalnya kemampuan menyusun tabel kebenaran, menerjemahkan syarat ke predikat, atau menyusun langkah bukti) yang ingin dicapai di akhir semester, selain nilai huruf.
  \item \textbf{Refleksi tahapan kurikulum:} Berdasarkan skema tahapan Bab 3--14 (Gambar~\ref{fig:peta-konsep}), pilih satu tahap yang menurut estimasi awal paling menantang; jelaskan alasan singkatnya.
  \item \textbf{Tokoh dan perkembangan logika:} Cari dua profil tokoh: (a) kontributor klasik pada logika formal, (b) kontributor pada aljabar Boolean atau fondasi logika dalam komputasi. Jelaskan secara singkat relevansi pemikiran mereka terhadap perangkat komputasi modern.
\end{aktivitas}

% ============================================================
% LATIHAN DAN REFLEKSI
% ============================================================
\begin{latihan}
  \item Jelaskan risiko \textit{logic bug} pada sistem perbankan bergerak jika penalaran kondisional (misalnya hubungan AND/OR) dan verifikasi melalui tabel kebenaran tidak diperlakukan secara sistematis.
  \item Bandingkan indikator keberhasilan berbasis luaran (OBE) dengan indikator berbasis kegiatan masukan (\textit{input-based}). Apa bukti konkret bahwa seorang mahasiswa telah mencapai capaian yang dimaksud dalam buku ini?
  \item Mengapa urutan pembelajaran bab-bab logika disarankan mengikuti silabus? Jelaskan dengan konsep ketergantungan konsep (\textit{prasyarat}) atau kurikulum spiral.
  \item Bahas peran logika predikat orde pertama dalam merumuskan aturan dan inferensi pada sistem berbasis pengetahuan; bandingkan secara singkat dengan pendekatan pembelajaran mesin yang tidak eksplisit memakai aturan logika.
\end{latihan}

% ============================================================
% ASESMEN
% ============================================================
\begin{asesmen}
\textbf{Instrumen pemahaman orientasi buku}

\textbf{A. Kuis pilihan ganda}

\begin{enumerate}
  \item Dalam pendekatan \textit{Outcome-Based Education} (OBE), penilaian utama pada mata kuliah komputasi selaras dengan prinsip berikut:
  \begin{enumerate}
    \item Tebalnya buku catatan di akhir semester.
    \item Demonstrasi kemampuan menerjemahkan pernyataan bahasa alami ke formalisme logika atau menyusun solusi yang memenuhi capaian.
    \item Kelengkapan absensi perkuliahan sebagai satu-satunya indikator.
    \item Hafalan nomor bab tanpa demonstrasi keterampilan.
  \end{enumerate}
  Jawaban: (b)

  \item Menurut pemetaan Bab 3--14, \textbf{tahap IV} (Bab 12--14) mengarah pada kompetensi berikut:
  \begin{enumerate}
    \item Sejarah filsafat logika klasik semata
    \item Sengaja memperkenalkan \textit{logic bug} pada rangkaian
    \item Teknik pembuktian, induksi matematika, serta argumen formal untuk validasi algoritma
    \item Hanya operasi tabel kebenaran dua variabel
  \end{enumerate}
  Jawaban: (c)
\end{enumerate}

\textbf{B. Esai singkat}

\begin{enumerate}
  \item Jelaskan rencana belajar mandiri untuk bagian \textit{Latihan} dan tabel kebenaran: bagaimana cara memverifikasi pemahaman tanpa bergantung pada menghafal solusi?
\end{enumerate}

\textbf{Rubrik penilaian}: Diatur dalam kontrak perkuliahan atau panduan asesmen dosen.
\end{asesmen}

% ============================================================
% CHECKLIST KOMPETENSI
% ============================================================
\begin{checklist}
  \item Saya mafhum dan tunduk pada urgensi membedah **Logika** sebagai modal \textit{Skill} terpenting pembeda Programmer Arsitek sejati lawan \textit{Coder} kacangan.
  \item Saya rela meninggalkan format pasif catatan bangku lama demi menyongsong mode demonstrasi maut ala spesifikasi standar industri berbasis luaran \textit{OBE}.
  \item Saya menguasai jalur rute mendaki Bab *(Roadmap)* buku ini, tahu dari mana ia lahir (Boolean dasar), hingga ke mana puncak mautnya (AI System Rules/SQL Predicates).
  \item Saya berikrar pantang membaca loncat bab, serta pantang percaya mencontek kunci tabel penulis sebelum merobek lembaran kertas mencoba menaklukkan masalah secara sepihak.
\end{checklist}

% ============================================================
% RANGKUMAN
% ============================================================
\begin{rangkuman}
Bab ini memperkenalkan struktur buku ajar, capaian berbasis OBE, serta hubungan logika formal dengan praktik teknik informatika. Orientasi ini menjadi acuan untuk mempelajari materi inti pada Bab 3--14 dan evaluasi integratif pada Bab 15.

\textbf{Poin kunci:}
\begin{itemize}
  \item \textbf{CPL, CPMK, Sub-CPMK:} capaian pembelajaran dijabarkan secara terukur; tabel pemetaan menghubungkan setiap Sub-CPMK dengan nomor bab.
  \item \textbf{OBE:} penekanan pada demonstrasi luaran (keterampilan dan bukti penguasaan), bukan hanya pada kegiatan masukan.
  \item \textbf{Peran logika:} mendukung analisis argumen, perancangan sistem yang dapat dijelaskan secara formal, dan pencegahan \textit{logic bug}.
  \item \textbf{Alur baca:} materi bersifat kumulatif; disarankan mengikuti urutan bab dan memanfaatkan aktivitas, latihan, serta checklist di setiap bab.
\end{itemize}

\textbf{Kata kunci:} orientasi buku, OBE, CPL, CPMK, Sub-CPMK, RPS, logika proposisional, logika predikat, logic bug, kurikulum spiral
\end{rangkuman}

\ifSubfilesClassLoaded{
  \renewcommand{\bibname}{Daftar Pustaka}
  \bibliographystyle{plain}
  \bibliography{../references}
}{}
\end{document}
